\section{The finite simple-group obstruction}\label{sec:finite-groups}

We prove Theorem~\ref{fin:obstruction}, following the finite-group
argument of \cite[Section~6]{MPNF}. In addition to the double-coset
obstruction proved there, we verify for each constructed set that its
complement in an abelian subgroup meeting it is a subgroup. In
characteristic two, this property supplies the finite abelian quotient
used in Proposition~\ref{rec:two}. The resulting functions have additive
period groups of finite index, which rules out the exceptional points
left by the density argument.

The proof uses the classification of finite simple groups
\cite[Classification Theorem, p.~737]{Aschbacher2004}. Its common
mechanism is to locate a part of a commuting pair on which centralizers
are small. For alternating and linear groups, this part consists of two
points or a line and a hyperplane. For the remaining classical groups it
is a large irreducible subspace. Most exceptional and sporadic groups
instead have an abelian subgroup whose nonidentity elements have exactly
that subgroup as centralizer. We first establish the two elementary tests
behind these constructions.

Throughout this section, $\mathcal F$ is a finite nonabelian simple group,
$C(B)=C_{\mathcal F}(B)$, and
\[
 I(\mathcal F)=\{g\in\mathcal F:g^2=1\},\qquad
 k(\mathcal F)=\text{the number of conjugacy classes of }\mathcal F.
\]
For commuting $W,Z\in\mathcal F$, the condition to be excluded is
\begin{equation}\label{fg:relation}
 ZW\in C(B)(ZW^{-1})C(BZ^2)
 \quad\text{for every }B\in\mathcal F.
\end{equation}
The choice of the nonempty proper conjugacy-invariant set $\mathcal S$
will be specified in each family. For every choice we must prove both
that Equation~\eqref{fg:relation} fails for some $B$ whenever
$W\in\mathcal S$ and $Z\notin\mathcal S$ commute, and that the
complement condition in Theorem~\ref{fin:obstruction}(ii) holds.

\subsection{Two elementary tests}

First take
\begin{equation}\label{fg:square-set}
 \mathcal S=\{W\in\mathcal F:W^2\ne1\}.
\end{equation}
This set is nonempty, since a group of exponent two is abelian, and it
does not contain the identity. For every abelian subgroup $A$, its
complement is the subgroup $A[2]=\{a\in A:a^2=1\}$. Thus this choice
automatically satisfies Theorem~\ref{fin:obstruction}(ii). For commuting $W\in\mathcal S$ and
$Z\notin\mathcal S$, put $v=ZW$. Then $Z^2=1$, $v^2=W^2\ne1$, and
Equation~\eqref{fg:relation} becomes
\begin{equation}\label{fg:inverse-coset}
 v\in C(B)v^{-1}C(B)\quad\text{for every }B\in\mathcal F.
\end{equation}
Suppose that $\mathcal F$ acts on a set and that some point $p$ satisfies
$v^2p\ne p$. If $C(B)$ fixes $p$ and $vp$ pointwise, writing
$v=a v^{-1}b$ with $a,b\in C(B)$ gives
$vp=a v^{-1}p$. Applying $a^{-1}$ gives $vp=v^{-1}p$, a contradiction.
We call this the point-pair test. We will also use its incidence version:
if $C(B)$ fixes a line $p$ and preserves a subspace $L$, while $vp\subset L$
and $v^{-1}p\not\subset L$, the same equation is impossible.

The second test uses a self-centralizing abelian subgroup.
\begin{lemma}\label{fg:abelian-test}
Suppose that $C\subset\mathcal F$ is an abelian subgroup of odd order
$c>1$ such that $C_{\mathcal F}(t)=C$ for every $1\ne t\in C$. If
\begin{equation}\label{fg:class-bound}
 |\mathcal F|>k(\mathcal F)c^4,
\end{equation}
then the union $\mathcal S$ of the conjugates of $C\setminus\{1\}$
satisfies both assertions of Theorem~\ref{fin:obstruction}.
\end{lemma}
\begin{proof}
The set is nonempty and excludes the identity. If $W\in\mathcal S$ and
$Z$ commutes with $W$, then $Z$ lies in the same conjugate of $C$ as $W$.
Consequently $Z\notin\mathcal S$ forces $Z=1$. More generally,
if an abelian subgroup $A$ meets $\mathcal S$, conjugating one member of
$A\cap\mathcal S$ into $C$ puts all of $A$ in $C$; hence
$A\setminus\mathcal S=\{1\}$. This proves the complement assertion.

Conjugate $W$ into $C$. Testing Equation~\eqref{fg:inverse-coset}
with every conjugate of a fixed nonidentity element of $C$ shows that
every conjugate $v$ of $W$ belongs to $Cv^{-1}C$. If
$v=a v^{-1}b$, $a,b\in C$, then
\[
 (v b^{-1})^2=a b^{-1}\in C.
\]
If this square is nonidentity, $v b^{-1}$ centralizes a nonidentity
element of $C$, so $v\in C$. Otherwise $v\in I(\mathcal F)C$.
Since $I(\mathcal F)$ is conjugacy-invariant, $I(\mathcal F)C=CI(\mathcal F)$.
Thus the whole conjugacy class of $W$ lies in $CI(\mathcal F)$, and
\[
 \frac{|\mathcal F|}{c}\le c|I(\mathcal F)|.
\]
The Frobenius--Schur formula and Cauchy--Schwarz give
\[
 |I(\mathcal F)|
 =\sum_{\chi\in\operatorname{Irr}(\mathcal F)}\nu_2(\chi)\chi(1)
 \le\sum_\chi\chi(1)
 \le\sqrt{k(\mathcal F)|\mathcal F|}.
\]
These inequalities contradict Equation~\eqref{fg:class-bound}.
\end{proof}

\subsection{Alternating, linear, and sporadic groups}

For alternating and projective special linear groups we use
Equation~\eqref{fg:square-set}. In $A_n$, choose $p$ with
$v^2p\ne p$. There is an even permutation $B$ fixing exactly $p,vp$
and having distinct odd cycle lengths greater than one on the remaining
points, except when $n=6$ or $8$. Indeed, for odd $n$ use the single
cycle of length $n-2$, and for even $n\ge10$ use lengths $3,n-5$.
Every permutation centralizing these moving cycles is even on their
support. Therefore a centralizer element in $A_n$ cannot interchange the
two fixed points. The point-pair test applies. The two exceptions are
$A_6\simeq\PSL_2(9)$ and $A_8\simeq\PSL_4(2)$.

Now let $\mathcal F=\PSL_l(q)$, $l\ge3$. Choose a line $p$ such that
$v^2p\ne p$. The three lines $p,vp,v^{-1}p$ are distinct. A hyperplane
$L$ can be chosen to contain $vp$ and contain neither $p$ nor $v^{-1}p$:
in the quotient by $vp$, choose a linear functional nonzero on both
remaining nonzero vectors. The union of two proper hyperplanes in the
dual is not the entire dual, also over $\bbF_2$.

On the direct sum $p\oplus L$, take $B$ to be the identity on $p$ and
an irreducible norm-one field multiplication on $L$. Such an element
exists in the cyclic group of order $(q^{l-1}-1)/(q-1)$: a generator
cannot lie in a proper subfield, since that group's order exceeds
$q^{(l-1)/2}-1$ when $l-1\ge2$. A projective centralizer of $B$ is an
actual centralizer. Indeed, scalar multiplication must preserve its
unique base-field eigenvalue, which is $1$, and hence the multiplier is
$1$. The centralizer fixes $p$ and preserves $L$, contradicting the
incidence version of Equation~\eqref{fg:inverse-coset}.

For simple $\PSL_2(q)$, the pointwise stabilizer of two points of the
projective line is a split torus. It contains an element whose
centralizer fixes both points, unless $q=5$: its order is $q-1$ for
even $q$ and $(q-1)/2$ for odd $q$, and a noninvolution in it is regular
with that torus as centralizer. The case $q=5$ is $A_5$. Choosing the
torus for $p,vp$ proves the assertion in all these cases.

For the sporadic groups other than $M_{12}$, and also for the Tits group,
Table~\ref{fg:table-sporadic} gives a prime $c$ for which a cyclic
subgroup of order $c$ is the centralizer of each of its nonidentity
elements. The centralizer orders, class numbers, and group orders are
those in the \emph{Atlas} and the online \emph{ATLAS of Finite Group
Representations}~\cite{Atlas,AtlasOnline}. In each row, an element $x$
in the class labeled $c\mathrm A$ in that group's conjugacy-class table
has $C(x)=\langle x\rangle$ of order $c$. Since $c$ is prime, every
nonidentity power of $x$ has the same centralizer.
The displayed $k$ is the exact class number. In every row the group
order exceeds $kc^4$, so Lemma~\ref{fg:abelian-test} applies.
For example, the smallest margin is $10\cdot5^4=6250<7920=|M_{11}|$.
One may also check the other inequalities using $k\le200$, except that
$k\le25$ suffices for $M_{22},M_{23},J_1$.

\begin{table}[htbp]
\centering
\small
\begin{tabular}{lrr@{\qquad}lrr}
\toprule
Group & $c$ & $k$ & Group & $c$ & $k$\\
\midrule
$M_{11}$ & 5 & 10 & $\mathrm{O'N}$ & 31 & 30\\
$M_{22}$ & 11 & 12 & $\mathrm{Co}_1$ & 23 & 101\\
$M_{23}$ & 23 & 17 & $\mathrm{Co}_2$ & 23 & 60\\
$M_{24}$ & 23 & 26 & $\mathrm{Co}_3$ & 23 & 42\\
$J_1$ & 7 & 15 & $\mathrm{Fi}_{22}$ & 13 & 65\\
$J_2$ & 7 & 21 & $\mathrm{Fi}_{23}$ & 23 & 98\\
$J_3$ & 19 & 21 & $\mathrm{Fi}_{24}'$ & 29 & 108\\
$J_4$ & 43 & 62 & $\mathrm{HN}$ & 19 & 54\\
$\mathrm{HS}$ & 11 & 24 & $\mathrm{Ly}$ & 67 & 53\\
$\mathrm{McL}$ & 11 & 24 & $\mathrm{Th}$ & 31 & 48\\
$\mathrm{He}$ & 17 & 33 & $\mathbb B$ & 47 & 184\\
$\mathrm{Ru}$ & 29 & 36 & $\mathbb M$ & 71 & 194\\
$\mathrm{Suz}$ & 13 & 43 & ${}^2F_4(2)'$ & 13 & 22\\
\bottomrule
\end{tabular}
\caption{Self-centralizing prime-order subgroups for the sporadic and
Tits groups.}\label{fg:table-sporadic}
\end{table}

For $M_{12}$ use its sharply $5$-transitive action of degree $12$ and
the \emph{Atlas} class $8B$, of cycle shape $1^2\,2\,8$
\cite{Atlas,AtlasOnline}; the online database labels this permutation
representation $12a$ (the other degree-$12$ action interchanges the
classes $8A$ and $8B$).
An element centralizing $B\in8B$ restricts to a cyclic shift on its
eight-point orbit. This restriction is injective: its kernel fixes
eight points, and an element fixing five points is the identity.
The centralizer is therefore exactly $\langle B\rangle$, and fixes
the two fixed points individually. Two-transitivity moves this pair to
$p,vp$, so the point-pair test applies with
Equation~\eqref{fg:square-set}.

\subsection{An auxiliary unitary count}

The remaining classical groups need a torus argument on a large
irreducible block. The relevant obstruction takes place in the full
unitary isometry group, which need not itself be simple.

\Needspace{9\baselineskip}
\begin{lemma}\label{fg:unitary-count}
Let $p\ge3$ be an odd prime, let $H=\U_p(q)$ be the full isometry
group of a nondegenerate Hermitian form over $\bbF_{q^2}$, and let
$T\subset H$ be an irreducible torus of order $q^p+1$. If $w\in T$
has field degree $p$ over $\bbF_{q^2}$, then
\[
 w\in T^h w^{-1}T^h\quad\text{for every }h\in H
\]
is impossible.
\end{lemma}
\begin{proof}
Put
\[
 \overline H=H/Z(H),\qquad \overline T=T/Z(H),\qquad
 c=\frac{q^p+1}{q+1},\qquad d=(p,q+1).
\]
The actual centralizer of $w$ is $T$. If $g$ centralizes its projective
image, then $gwg^{-1}=\lambda w$ for a scalar $\lambda$ of norm one.
Taking determinants gives $\lambda^p=1$, so
\begin{equation}\label{fg:unitary-centralizer}
 |C_{\overline H}(\overline w)|\le dc.
\end{equation}

Suppose the displayed condition in the lemma holds. Every conjugate
$v$ of $w$ then lies in $Tv^{-1}T$. The calculation used in
Lemma~\ref{fg:abelian-test} gives an element $u=v b^{-1}$ whose
square lies in $T$. If $u^2$ is nonscalar, it has field degree $p$,
because $p$ is prime; its centralizer is $T$, and hence $u\in T$.
If $u^2$ is scalar, $\overline u$ has square one. Consequently the
entire projective conjugacy class of $\overline w$ lies in
$\overline T I(\overline H)$.

The determinant quotient of $\overline H$ has order $d$, and
$\overline T$ maps onto it. To see this, write $T$ as the norm-one
subgroup of $\bbF_{q^{2p}}^\times$ over $\bbF_{q^p}$. The determinant
of a field multiplication is its norm to $\bbF_{q^2}$. On this torus,
that norm takes a generator to a generator of the scalar norm-one
group of order $q+1$; its exponent is congruent to $c$ modulo
$q^p+1$. Modding out by scalar matrices mods the determinant out by
$p$th powers, giving the quotient of order $d$.

All members of the conjugacy class have the same determinant in this
quotient. For each fixed $i\in I(\overline H)$, at most $c/d$ elements
$t\in\overline T$ have $ti$ with that determinant. Combining this
with Equation~\eqref{fg:unitary-centralizer} gives
\[
 \frac{|\overline H|}{dc}
 \le \frac c d |I(\overline H)|,
 \qquad\text{and therefore}\qquad
 |\overline H|\le c^2|I(\overline H)|.
\]
We show that the reverse strict inequality always holds.

Every projective involution lifts to an involution in $H$. Indeed,
if $g^2=\alpha I$, then $\alpha^p=\det(g)^2$ in the scalar cyclic
group of order $q+1$. As $p$ is odd, $\alpha$ is a square in that
group, and scalar rescaling gives an involution.

For odd $q$, an involution is determined up to conjugacy by the
dimensions $a,p-a$ of its two nondegenerate eigenspaces. Projectively
the unordered partitions are indexed by $1\le a<p/2$, and their
centralizers have order at least
\begin{equation}\label{fg:unitary-odd-centralizer}
 \frac{|\U_a(q)|\,|\U_{p-a}(q)|}{q+1}.
\end{equation}
For even $q$, the Jordan type is $2^a1^b$, where $b=p-2a$. Its full
unitary centralizer has order
\begin{equation}\label{fg:unitary-even-centralizer}
 q^{a^2+2ab}|\U_a(q)|\,|\U_b(q)|.
\end{equation}
We recall the structure behind this formula, including in characteristic
two; see also \cite[Section~2.6, Case~(A)(ii),(iv), pp.~34--36]{Wall1963}.
In the ambient general
linear group the connected matrix centralizer has reductive quotient
$\GL_a\times\GL_b$ and unipotent radical of dimension $a^2+2ab$.
The Hermitian Frobenius gives the unitary forms on both multiplicity
spaces. Equivalently, for $g=1+N$ of order two, $N=N^*$; the induced
forms on $\ker N/\im N$ and, using $N$, on the quotient by $\ker N$
are the forms on these two spaces. Nondegenerate Hermitian forms of
a given dimension over a finite field are isometric. Connectedness
and Lang's theorem give one full-unitary class for each occurring
type and the stated order. Dividing by $q+1$ gives a lower bound
for its projective centralizer.

The unitary order formula implies
\begin{equation}\label{fg:unitary-order-lower}
 |\U_j(q)|=q^{j^2}\prod_{r=1}^j(1-(-q)^{-r})\ge q^{j^2}.
\end{equation}
For completeness, pair consecutive factors: their product is
\[
 (1+q^{-(2r-1)})(1-q^{-2r})
 =1+q^{-2r}(q-1-q^{-(2r-1)})>1.
\]
An unpaired final factor also exceeds one. Thus every nonidentity
projective involution class has centralizer order at least
\[
 M=\frac{q^{(p^2+1)/2}}{q+1}.
\]
There are at most $(p-1)/2$ such classes, and including the identity
gives
\[
 \frac{|I(\overline H)|}{|\overline H|}\le\frac{p+1}{2M}.
\]
When $p\ge5$, we have $c<q^{p-1}$ and
\[
 \frac{M}{q^{2p-2}}
 =\frac{q^{(p^2-4p+5)/2}}{q+1}
 \ge\frac{2^p}{3}>\frac{p+1}{2}.
\]
The first inequality uses $(p^2-4p+5)/2\ge p$ and monotonicity in
$q\ge2$. This proves $c^2|I(\overline H)|<|\overline H|$.

It remains to check $p=3$. Now
\[
 c=q^2-q+1,\qquad h=|\overline H|=q^3(q^3+1)(q^2-1).
\]
There is one nonidentity involution class. In odd characteristic its
centralizer lower bound is $M=q(q-1)(q+1)^2$, and $h/M=q^2c$.
In even characteristic it is $M=q^3(q+1)$, and $h/M=(q^2-1)c$.
In both cases $|I(\overline H)|\le1+q^2c$, whereas
\[
 h-c^2(1+q^2c)
 =c\bigl((3q-4)q^4+(q-2)q^2+q-1\bigr)>0
 \quad(q\ge2).
\]
This includes $\U_3(2)$ and finishes the proof.
\end{proof}

\subsection{Symplectic and orthogonal groups}

Let $\mathcal F$ be a projective symplectic or orthogonal simple group
on its natural space of dimension $2l$ or $2l+1$. We first treat
$l\ge3$, using the standard low-rank identifications and treating
$\PSp_4$ separately below. Choose an odd prime
\begin{equation}\label{fg:block-prime}
 l/2<p\le l,
\end{equation}
requiring $p<l$ in plus orthogonal type. For $l\ge4$, set
$m=\lfloor l/2\rfloor\ge2$. Bertrand's postulate gives
$m<p<2m$, hence $l/2<p<l$, and this prime is odd.
For $l=3$ take $p=3$, apart from plus type, which is $\PSL_4$ and
has already been treated. In odd-dimensional orthogonal type we may
assume odd characteristic, since in even characteristic the group
has the corresponding symplectic realization.

We use a nondegenerate block $E$ of dimension $2p$ carrying a torus
of order $q^p+1$. In orthogonal type this block has minus type; the
complement can be chosen with the sign needed for the ambient form.
For minus type $p=l$ the block is the whole space, whereas the
requirement $p<l$ in plus type leaves a complement of the necessary
sign. The torus is the norm-one subgroup of
$\bbF_{q^{2p}}^\times$ over $\bbF_{q^p}$, acting on $E$ by field
multiplication.

Define $\mathcal S$ to consist of those projective elements whose
square, on the natural space, has an irreducible self-dual block of
dimension $2p$. This condition is independent of the choice of a
scalar lift and is conjugacy-invariant. The block is unique, because
$4p>2l+1$, and is nondegenerate: its dual would otherwise require a
second block of the same degree. The set excludes identity and is
nonempty. Indeed, if $t$ generates the torus of order $q^p+1$, use
$t^2$ on the block and identity on the complement. Squaring meets
the orthogonal component and spinor conditions, whose quotients
have exponent two. Moreover $t^4$ still has degree $2p$. To see
this, the quotient of the norm-one torus by its intersection with
$\bbF_{q^2}^{\times}$ has odd order
\[
 \frac{q^p+1}{q+1}>1.
\]
Thus $t^4\notin\bbF_{q^2}$. The only other maximal proper subfield
is $\bbF_{q^p}$, where a norm-one element is $\pm1$ and hence already
lies in $\bbF_{q^2}$.

Suppose now that $W\in\mathcal S$, that $Z\notin\mathcal S$
commutes with $W$, and that Equation~\eqref{fg:relation} holds.
Lift $W,Z$ to natural isometries, using $\Omega$ in orthogonal type.
Their possible projective commutation multiplier has order at most
two. They therefore commute actually with $W^2$ and preserve its
unique marked block $E$. On $E$ they act as field multiplications
$w,z\in\bbF_{q^{2p}}^\times$ of norm one. Since $Z\notin\mathcal S$,
$z^2$ lies in a proper subfield. The norm-one condition puts $z^2$
in $\bbF_{q^2}$, and then $z$ itself is in $\bbF_{q^2}$: its degree
over that field divides both $p$ and $2$.

There is a full unitary group $H=\U_p(q)$ acting on $E$ and
containing its field torus $T$. One way to see this directly is to
write the invariant form in field coordinates. Its bilinear form
has the shape
\[
 \Tr_{\bbF_{q^{2p}}/\bbF_q}(\kappa x y^*),\qquad x^*=x^{q^p}.
\]
Take the trace first to $\bbF_{q^2}$. If the resulting form is
anti-Hermitian and the characteristic is odd, multiply it by
$\delta\in\bbF_{q^2}^{\times}$ with $\delta^q=-\delta$; this makes
it Hermitian without changing its isometry group. In characteristic
two an anti-Hermitian form is already Hermitian. This gives the
required unitary structure. In the
quadratic case the form is
\[
 Q(x)=\Tr_{\bbF_{q^p}/\bbF_q}(\kappa xx^*),
 \qquad\kappa\in\bbF_{q^p}^\times.
\]
In characteristic two this expression follows from the polarization
and $W^2$-invariance: two invariant quadratic forms with the same
polarization differ by the square of an invariant linear form,
and an irreducible block of degree greater than one has no such
nonzero form. More explicitly, the Hermitian form
$h(x,y)=\Tr_{\bbF_{q^{2p}}/\bbF_{q^2}}(\kappa xy^*)$
in the quadratic case satisfies $h(x,x)=Q(x)$: on $\bbF_{q^p}$ the
trace to $\bbF_{q^2}$ equals the trace to $\bbF_q$, since $p$ is odd.
Thus $H$ preserves the original form on $E$ in both cases. The element
$z$ is a scalar of this unitary group.

For each $h\in H$, choose $B$ acting on $E$ by a square of a
generator of $T^h$ and as identity on the complement. This is an
allowed ambient element. Even if $h$ is not in $\Omega$, conjugation
by an isometry preserves $\Omega$, so the orthogonal component
condition is unchanged. Both $B$ and $BZ^2$ have a unique marked
square block $E$: multiplication of a full-degree element by a
base-field scalar preserves its degree over $\bbF_{q^2}$, and the
norm-one condition then gives full degree $2p$ over $\bbF_q$.

Each projective centralizer of $B$ or $BZ^2$ commutes actually with
its square, preserves $E$, and restricts there into $T^h$. This
does not require the action on the complement to be semisimple;
the complement is too small to have an irreducible factor of
degree $2p$. Restricting Equation~\eqref{fg:relation} to $E$
and absorbing its scalar multipliers and the scalar $z$ into $T^h$
therefore gives
\[
 w\in T^h w^{-1}T^h\quad\text{for every }h\in H.
\]
The element $w$ has degree $p$ over $\bbF_{q^2}$. This contradicts
Lemma~\ref{fg:unitary-count}.

To verify the complement assertion, let $A$ be abelian and choose
$W\in A\cap\mathcal S$. Its marked block $E$ is now fixed.
The preceding centralizer argument gives a well-defined homomorphism
\begin{equation}\label{fg:classical-restriction}
 \rho:A\longrightarrow
 T/(T\cap\bbF_{q^2}^{\times}).
\end{equation}
Indeed every lift restricts on $E$ to a field multiplication, and
changing the lift multiplies it by an ambient central scalar lying in
$\bbF_{q^2}^{\times}$. The complement of $E$ is too small to carry a
marked block. For a restriction $z\in T$, its square has proper degree
over $\bbF_q$ precisely when $z^2\in\bbF_{q^2}$: the other maximal
proper subfield is $\bbF_{q^p}$, whose norm-one elements are $\pm1$.
Because $p$ is odd, $z^2\in\bbF_{q^2}$ is equivalent to
$z\in\bbF_{q^2}$. Full degree automatically gives self-duality by the
norm-one equation. Thus $A\setminus\mathcal S=\ker\rho$, as required.

\subsection{Projective unitary groups and rank two symplectic groups}

Let $\mathcal F=\PSU_l(q)$, $l\ge3$, excluding the nonsimple pair
$(l,q)=(3,2)$. Choose an odd prime $l/2<p\le l$. Define
$\mathcal S$ by requiring the square of a natural lift to have an
irreducible block of degree $p$ over $\bbF_{q^2}$. As above this
block is unique and nondegenerate. If $l>p$, the determinant of a
chosen torus element on the block can be corrected on its nonzero
dimensional orthogonal complement. If $l=p$, use instead the
determinant-one torus of order
\[
 c=\frac{q^p+1}{q+1}.
\]
Its scalar intersection has order $d=(p,q+1)$. Except for
$(p,q)=(3,2)$, $c>d$: for $p\ge5$,
$c\ge q^{p-2}(q-1)+1\ge2^{p-2}+1>p$, and for $p=3,q\ge3$,
$c\ge7>3$. As $c$ is odd, a generator and its square are both
nonscalar and hence have degree $p$. This proves that
$\mathcal S$ is nonempty; it again excludes identity.

We check the projective-centralizer issue before applying the block
argument. Suppose $gAg^{-1}=\lambda A$, where $A\in\SU_l(q)$ has
the marked square block. Multiplication by a base-field scalar
preserves irreducible degrees, so $g$ preserves the unique block.
The determinant on the block gives $\lambda^p=1$, and the total
determinant gives $\lambda^l=1$. If $l>p$, then $p<l<2p$, so these
conditions force $\lambda=1$. If $l=p$ and $\lambda\ne1$, multiplication
by $\lambda$ cycles all $p$ distinct eigenvalues. Writing $w$ for any one of them,
their product is $w^p$, since $p$ is odd, and determinant one gives
$w^p=1$.
But $\lambda$ has order $p$ in the scalar group, so $p\mid q+1$;
all $p$th roots of unity then lie in $\bbF_{q^2}$, contradicting
the full degree of $w$. Thus these projective centralizers are
actual centralizers on the block.

For commuting $W\in\mathcal S$ and $Z\notin\mathcal S$, the
restrictions are consequently field scalars $w,z$, with
$z\in\bbF_{q^2}$ by the same odd-degree argument as before.
For every conjugate $T^h$ in the full block unitary group choose
$B$ with irreducible square on that torus, correcting its determinant
on the complement if necessary. The square of $BZ^2$ retains full
degree. Applying the preceding projective-centralizer argument
to $B$ and $BZ^2$, Equation~\eqref{fg:relation} restricts to
the relation prohibited by Lemma~\ref{fg:unitary-count}.

For the complement assertion, fix $W\in A\cap\mathcal S$ in an
abelian subgroup $A$. The projective-centralizer argument just proved
gives a restriction homomorphism as in
Equation~\eqref{fg:classical-restriction}, now with $T$ the torus on the
$p$-dimensional Hermitian block. A restriction $z$ has square of degree
less than $p$ precisely when $z^2$, equivalently $z$, is in
$\bbF_{q^2}$. Since the complementary subspace has dimension less than
$p$, this characterizes nonmembership in $\mathcal S$. Once again
$A\setminus\mathcal S=\ker\rho$.

For $\PSp_4(q)$ take $\mathcal S$ to be the set whose square is
irreducible on the natural four-dimensional space. It is nonempty
by the torus of order $q^2+1$ and excludes identity. Lifting a
commuting pair $W\in\mathcal S$, $Z\notin\mathcal S$ gives norm-one
field scalars $w,z\in\bbF_{q^4}^\times$. Since $z^2\in\bbF_{q^2}$,
the norm-one condition gives $z^2=\pm1$. Thus $Z^2=1$ projectively,
and $v=ZW$ still has an irreducible square.

Choose a nonzero vector $x$ such that $x,vx$ span an isotropic
plane $L$. Such an $x$ exists in the field description of the
alternating form: the condition
$\Tr_{\bbF_{q^4}/\bbF_q}(\kappa x(vx)^*)=0$ is one homogeneous
$\bbF_q$-linear condition on $xx^*\in\bbF_{q^2}$. It has a nonzero
solution, and the field norm $x\mapsto xx^*$ is surjective.
Irreducibility ensures that $x,vx$ are independent and that
$v^{-1}x\notin L$; otherwise $v$ would satisfy a quadratic polynomial.

Choose a regular unipotent $B$ whose unique line-plane flag is
$\langle x\rangle\subset L$. Its existence in every characteristic
can be seen in a symplectic flag basis from the matrix
\[
 \begin{pmatrix}
 1&1&0&0\\0&1&1&-1\\0&0&1&-1\\0&0&0&1
 \end{pmatrix},
\]
for the alternating form with nonzero upper anti-diagonal entries
$1,1$. This matrix preserves the form and has one Jordan block,
also in characteristic two. Its centralizer preserves
$\ker(B-1)=\langle x\rangle$ and $\ker((B-1)^2)=L$.
A projective centralizer is actual, since a scalar multiple of
a unipotent matrix can have the same eigenvalues only when the
scalar is one. The incidence test now contradicts
Equation~\eqref{fg:inverse-coset}. If an abelian $A$ meets this $\mathcal S$, fix the irreducible block
using any member of the intersection. Every other member restricts to a
norm-one field scalar $z$, and the calculation above gives
\[
 z^2\text{ has proper degree}
 \quad\Longleftrightarrow\quad z^2\in\bbF_{q^2}
 \quad\Longleftrightarrow\quad z^2=\pm1.
\]
The last condition says precisely that its projective image has square
one. Thus $A\setminus\mathcal S=A[2]$. The nonsimple group
$\Sp_4(2)$ is not needed; its simple derived group is $A_6$.

\subsection{Small tori in exceptional groups}

We have proved the obstruction for all classical simple groups.
For most exceptional families we apply
Lemma~\ref{fg:abelian-test} to a rational maximal torus.
We use the usual simply connected algebraic realization followed
by its central quotient. Rational tori are obtained by Weyl-group
twisting, and their orders are determinants of the Frobenius action
minus one on their character lattices; see \cite{Steinberg} and
\cite[Propositions~25.1--25.3, pp.~219--220]{MalleTesterman}. Semisimple centralizers
in a simply connected semisimple group are connected
\cite[Remark~2.10]{SteinbergRegular1965}. The tori below and their
self-centralizing property in the simple quotient occur in
\cite[Corollary~8.2 and Lemma~8.7]{SegevSeitz2002}.
We verify that property here by a root-lattice estimate before applying
Lemma~\ref{fg:abelian-test}. Their orders outside type $G_2$
are also recorded in \cite[Section~3]{GarionLarsenLubotzky}; the two
$G_2$ orders follow directly from its order-three and order-six Weyl
actions.

Write $\Phi_j$ for the $j$th cyclotomic polynomial. Use the tori
in Table~\ref{fg:table-exceptional-tori}, before passing to the
central quotient. For $G_2(q)$ choose the sign avoiding a factor
$3$, and for $q=3$ choose the smaller order $7$.
\begin{table}[htbp]
\centering
\begin{tabular}{lll}
\toprule
Group & Torus order & Weyl action\\
\midrule
$G_2(q)$, $q\ge3$ & $q^2\pm q+1$ & Order $3$ or $6$\\
${}^3D_4(q),F_4(q)$ & $\Phi_{12}(q)$ & Characteristic polynomial $\Phi_{12}$\\
${}^{\epsilon}E_6(q)$ & $\Phi_9(\epsilon q)$ & $\Phi_9$, with diagram sign\\
$E_8(q)$ & $\Phi_{30}(q)$ & Coxeter action\\
\bottomrule
\end{tabular}
\caption{Tori used for the exceptional simple groups;
$\epsilon=1$ denotes split $E_6$ and $\epsilon=-1$ twisted $E_6$.}
\label{fg:table-exceptional-tori}
\end{table}

These actions are irreducible over $\bbQ$. For $G_2,F_4,E_8$ use
a Coxeter element or an appropriate power. In trialitarian $D_4$,
let $b$ be the central node, $a$ an outer node, and $\gamma$ the
diagram cycle of the outer nodes. The action $s_a s_b\gamma$ has
characteristic polynomial $X^4-X^2+1$. For $E_6$, the regular
order-nine Weyl element in Springer's
tables~\cite[Section 5.4, Table 1]{Springer} has
a primitive ninth-root eigenvalue. Its rational characteristic
polynomial, of degree six, is therefore $\Phi_9$. Its negative
belongs to the nontrivial diagram coset and gives the twisted
case.

After passing to the central quotient, dividing by $(3,q-\epsilon)$
in type $E_6$, every prime divisor of these torus orders is a
primitive cyclotomic prime. Their respective lower bounds are
\[
 7,\qquad13,\qquad19,\qquad31.
\]
Indeed a prime dividing $\Phi_m(q)$ either has multiplicative
order $m$ for $q$ modulo that prime, or is an exceptional prime
dividing $m$. The displayed polynomials have no such exceptional
factor except for $3$ in $\Phi_9(\epsilon q)$; when present this
factor has valuation one and is exactly removed by the center.
All resulting torus orders are odd and greater than one.

We next prove that \emph{every} nonidentity element of each resulting
torus has centralizer exactly that torus. It suffices to do so for
an element $t$ of prime order $r$ upstairs, with $r$ one of these
primitive primes. If a root $\alpha$ vanished on $t$, Frobenius
invariance would make every root in its twisted Weyl orbit vanish
on $t$. By irreducibility that orbit spans the rational character
space. Choose a linearly independent subset. The lattice it
generates has full rank in the weight lattice, and its index must
be divisible by $r$, since evaluation at $t$ is a nontrivial
character of order $r$ trivial on this sublattice.

Normalize long roots to have squared length two. Hadamard's
inequality bounds these lattice indices by
\begin{equation}\label{fg:root-index}
 2\sqrt3,\qquad 8,\qquad 8\sqrt3,\qquad16
\end{equation}
in the respective rows of Table~\ref{fg:table-exceptional-tori}.
The respective weight-lattice covolumes are
\[
 1/\sqrt3,\qquad1/2,\qquad1/\sqrt3,\qquad1.
\]
The second value holds for both $D_4$ and $F_4$.
Each bound is smaller than the corresponding $r$.
Thus $t$ is regular. Connectedness of its algebraic centralizer
makes that centralizer exactly the torus.

Any nonidentity element of a torus in the simple quotient has a
nontrivial prime-order power of this kind. The argument survives
projective centralization: one can lift that power with order $r$
coprime to the center, and a conjugate differing from it by a
central element must have that central factor equal to one,
by preservation of its order. Therefore its projective centralizer
is precisely the torus in the quotient. The hypotheses on $C$
in Lemma~\ref{fg:abelian-test} are now verified.

For the numerical inequality, we use
\begin{equation}\label{fg:exceptional-class-bound}
 k(\mathcal F)\le100q^l,
\end{equation}
where $l$ is the absolute rank. This follows, with a smaller
constant, from \cite[Theorem~1.1, p.~3024]{FulmanGuralnick} for the simply
connected fixed-point groups; taking a quotient cannot increase
the number of conjugacy classes. The standard group order
formulas \cite[Tables~24.1--24.2, pp.~208, 211]{MalleTesterman} give
\[
 \begin{array}{c|ccccc}
 \mathcal F&G_2(q)&{}^3D_4(q)&F_4(q)&{}^{\epsilon}E_6(q)&E_8(q)\\
 \hline
 |\mathcal F|>&0.7q^{14}&0.7q^{28}&q^{52}/2&q^{78}/6&q^{248}/2.
 \end{array}
\]
In the two small $G_2$ cases the direct comparisons are
\[
 \begin{aligned}
 |G_2(3)|&=4245696>100\cdot3^2\cdot7^4,\\
 |G_2(4)|&=251596800>100\cdot4^2\cdot13^4.
 \end{aligned}
\]
For $q\ge5$ the chosen $G_2$ torus has order $c\le1.24q^2$, and
$0.7q^{14}>100q^2(1.24q^2)^4$ already at $q=5$, with increasing
ratio thereafter. In trialitarian type $c<q^4$, so it is enough
that $0.7q^8>100$, true at $q=2$. For $F_4,E_6,E_8$ use respectively
$c<q^4,c<2q^6,c<2q^8$; the displayed lower bounds exceed
$100q^l c^4$ already at $q=2$, and thereafter with increasing
ratio. This proves the desired inequality in every row, so
Lemma~\ref{fg:abelian-test} completes these families.

\subsection{Type \texorpdfstring{$E_7$}{E7} and the Suzuki--Ree groups}

For $E_7$ we transfer the torus obstruction from a suitable $E_6$
subgroup. In the simply connected algebraic group choose a
maximal-rank subgroup $M$ geometrically a Levi subgroup with
derived group $D_M=[M,M]$ of type $E_6$ and one-dimensional center. Both
rational forms $E_6$ and ${}^2E_6$ can be obtained: the ambient
Weyl group contains $-1$, which normalizes this subsystem and
induces its nontrivial diagram option, so rational Weyl twisting
gives the second form. Choose the sign $\epsilon$ such that
$(3,q-\epsilon)=1$.

Conjugation on the derived group gives a map $\pi$ from $M(q)$ to
adjoint ${}^{\epsilon}E_6(q)$. Its image is the simple group of
that type. Indeed the simply connected derived subgroup already
maps onto it: both the finite center and the Frobenius cohomology
of its order-three algebraic center vanish under the chosen
coprimality condition. Denote this finite simple image by
$\mathcal F_0$, and let $C_0\subset\mathcal F_0$ be the torus used
above. Geometrically $M$ is $(E_{6,\mathrm{sc}}\times\bbG_m)/\mu_3$;
the central one-dimensional torus has Frobenius action $\epsilon q$
on its character lattice. The kernel $K_M=\ker\pi$ is therefore
central, with order dividing a product of $q-\epsilon$ and powers of $3$. It has no
prime factor in common with $|C_0|$.
Since $\pi(D_M(q))=\mathcal F_0$, every element of $M(q)$ differs
from an element of $D_M(q)$ by an element of $K_M$. Thus
\[
 M(q)=D_M(q)K_M,\qquad
 M(q)/D_M(q)\simeq K_M/(K_M\cap D_M(q)).
\]
In particular, for every primitive prime $r$ dividing $|C_0|$,
this quotient has order prime to $r$, and every element of
$r$-power order in $M(q)$ belongs to $D_M(q)$.

Define $\mathcal S$ in projective $E_7(q)$ to be the union of the
conjugates of the images of those elements of $M(q)$ whose
projection belongs to $C_0\setminus\{1\}$. The preimage in $M$
of the algebraic torus containing $C_0$ is a maximal torus of
$M$, and also of $E_7$. Thus all these elements are semisimple.
The set is nonempty by surjectivity, and it excludes identity:
the ambient central elements project trivially to adjoint $E_6$.

We need a prime-order power of each such element to be regular not
only in $E_6$ but in $E_7$. For $W\in M(q)$ with nonidentity projection
in $C_0$, choose a prime $r$ dividing the order of $\pi(W)$ and a
power $t$ of $W$ of order $r$. Because $K_M$ has order prime to $r$,
$\pi(t)\ne1$; the quotient calculation above also puts $t$ in $D_M(q)$.
Every $E_7$ root restricts to a nonzero weight on this $E_6$. To check
nonvanishing, the fundamental weight perpendicular to the $E_6$
subsystem has squared norm $3/2$; a root proportional to it would
have rational proportionality factor of square $4/3$, impossible.
The restriction has length at most $\sqrt2$. Its orbit under
the irreducible twisted Weyl action spans the six-dimensional
weight space. The same index argument as in
Equation~\eqref{fg:root-index} bounds the resulting
weight-lattice index by $8\sqrt3<19$, smaller than the primitive
prime. Thus no root vanishes on $t$, and its algebraic centralizer in $E_7$
is precisely the maximal torus containing it.

Now conjugate $W\in\mathcal S$ into the specified torus in $M(q)$.
Every element commuting with it projectively belongs to this
torus: the center of simply connected $E_7$ has order at most
two, so the primitive-prime power is actually centralized, and
its connected centralizer is the torus. In particular a commuting
$Z\notin\mathcal S$ lies in $M(q)$ and projects to $1$ in
$\mathcal F_0$; a nonidentity projection would lie in $C_0$ and
put $Z$ in $\mathcal S$.

For each conjugate of a nonidentity element of $C_0$ in
$\mathcal F_0$, choose a lift $B$ in $M(q)$. Both $B$ and $BZ^2$
have that nonidentity projection, so the preceding regularity
argument applies to each. Their projective centralizers are
toral and project into the corresponding conjugate of $C_0$.
Projecting Equation~\eqref{fg:relation} to $\mathcal F_0$
therefore puts the nonidentity projection $w$ of $W$ in
\[
 C_0^h w^{-1}C_0^h\quad\text{for every }h\in\mathcal F_0.
\]
The counting proof of Lemma~\ref{fg:abelian-test}, already
verified for this $E_6$ group, excludes this. This proves the
obstruction in $E_7$.

We must also identify the complement inside an abelian subgroup $A$
meeting $\mathcal S$. Fixing a member of the intersection conjugates
$A$ into the image $\overline T$ of the torus just used. The map
$\pi:T\to C_0$ descends to $\overline T$, since the ambient center
acts trivially on adjoint $E_6$. Every element with nonidentity
projection belongs to $\mathcal S$ by definition. Conversely, every
element of $\mathcal S$ has order divisible by a prime dividing
$|C_0|$: its nonidentity projection has such order, and passage to
projective $E_7$ removes at most a factor two. No element of the image of
$K_M$ has such a prime divisor. Therefore even conjugacy into a different
torus cannot put an element of $\ker\pi$ into $\mathcal S$, and
\[
 A\setminus\mathcal S=A\cap\ker(\pi:\overline T\to C_0).
\]
This completes both assertions in type $E_7$.

For the simple Suzuki and rank-one Ree groups, use
Equation~\eqref{fg:square-set} and their doubly transitive
rank-one BN-pair action. A pair stabilizer has a torus of order
$q-1$, containing a prime-order element of order greater than
three. In Suzuki type $q=2^{2a+1}\ge8$, and $q-1$ is odd and not
divisible by three. In Ree type $q=3^{2a+1}\ge27$, the number
$q-1$ has 2-adic valuation one and an odd factor greater than one,
not divisible by three.

Such a prime-order element $t$ is regular. On the standard
Frobenius-stable pair torus, the exceptional Frobenius $F$ acts
on roots by $F^*\alpha=\ell^j\alpha'$, where $\ell$ is the
characteristic and $\alpha'$ is a root of the opposite length;
the factor $\ell^j$ includes the field-Frobenius part.
If $\alpha(t)=1$, then $F(t)=t$ gives
$\alpha'(t)^{\ell^j}=1$. Since the order of $t$ is prime to
$\ell$, also $\alpha'(t)=1$. The two roots are nonparallel,
because a reduced root system has no parallel roots of different
lengths. The index of the
lattice they span in the weight lattice is at most $2\sqrt2$
in $B_2$ and at most $2$ in $G_2$, by the same length and
covolume calculation as above. Neither index can be divisible
by the chosen prime. Its centralizer is therefore the pair torus,
which fixes both points individually. Choose the pair to be
$p,vp$ and apply the point-pair test.

Finally consider ${}^2F_4(q)$, $q=2^{2a+1}\ge8$. Let $\gamma$ be
the normalized diagram reversal in $F_4$. The twisted Weyl
action $s_1s_2\gamma$ has square $s_1s_2s_4s_3$, a Coxeter
element. The resulting finite torus $C$ is contained in the
Frobenius-square torus of order $\Phi_{12}(q)$. Its order is one
of the two factors
\[
 q^2\pm\sqrt{2q^3}+q\pm\sqrt{2q}+1,
\]
with the signs chosen together; these are Malle's two cyclic tori,
as recorded in \cite[Theorem~8]{GarionLarsenLubotzky}, and their
orders multiply to $\Phi_{12}(q)$. Equivalently, the
determinant formula bounds its order between
$(\sqrt q-1)^4$ and $4q^2$. It is odd and greater than one.
Every prime divisor is at least thirteen, and the $F_4$
root-lattice argument proves that every nonidentity element
has centralizer precisely $C$. The standard group order gives
$|{}^2F_4(q)|>q^{26}/2$. Also
$k({}^2F_4(q))\le100q^4$; the sharper bound
$q^2+4q+17$ is given in \cite[Table~1, p.~3049]{FulmanGuralnick}.
Thus
\[
 |{}^2F_4(q)|>100q^4(4q^2)^4\ge k({}^2F_4(q))|C|^4
 \qquad(q\ge8),
\]
and Lemma~\ref{fg:abelian-test} applies.

\begin{proof}[Completion of the proof of Theorem~\ref{fin:obstruction}]
The classification consists of the alternating groups, the classical
and exceptional groups of Lie type, and the sporadic groups. All
families have been covered above. The low-rank orthogonal groups
have their usual linear, unitary, or symplectic realizations;
in particular plus type $l=3$ is $\PSL_4$, and rank two symplectic
type was treated separately. The exceptional small derived simple
groups are $G_2(2)'\simeq\PSU_3(3)$,
${}^2G_2(3)'\simeq\PSL_2(8)$, and the Tits group, already included
in Table~\ref{fg:table-sporadic}. Every chosen $\mathcal S$ is
nonempty, proper, and conjugacy-invariant, and for it
Equation~\eqref{fg:relation} has been excluded for every
commuting pair with $W\in\mathcal S$, $Z\notin\mathcal S$. The
complement assertion was established with each construction: it is
$A[2]$ for the square sets, $\{1\}$ for the self-centralizing torus
sets, and the kernel of a restriction or projection homomorphism in the
remaining cases. Finally, if $s\in\mathcal S$ but
$s^{-1}\notin\mathcal S$, the complement subgroup in $\langle s\rangle$
would contain $s^{-1}$ and hence $s$, a contradiction. This proves
invariance under inversion and completes Theorem~\ref{fin:obstruction}.
\end{proof}
